\daytitle{11}{圆周运动}

\begin{learningbox}{本讲义大致内容}
圆周运动是曲线运动最重要的形式。今天学习描述圆周运动的物理量——线速度、角速度、周期，以及向心加速度和向心力。从水平圆周到竖直圆周，建立完整的圆周运动分析框架。
\end{learningbox}

\section*{一、学习目标}

\begin{enumerate}
  \item 掌握线速度 $v = \dfrac{2\pi r}{T}$ 和角速度 $\omega = \dfrac{2\pi}{T}$ 及关系 $v = \omega r$。
  \item 理解向心加速度 $a_n = \dfrac{v^2}{r} = \omega^2 r$，方向指向圆心。
  \item 掌握向心力 $F_n = m\dfrac{v^2}{r}$，会分析水平圆周运动的受力。
  \item 会分析竖直圆周运动最高点和最低点的受力。
  \item 理解匀速圆周运动是变加速曲线运动（加速度方向不断变化）。
\end{enumerate}

\section*{二、初高中衔接}

初中只简单学过圆周运动的“快慢”（用转速描述），没有量化公式。

高中引入三个核心物理量：
\begin{itemize}
  \item 线速度 $v$：沿圆周切线方向，$\mathrm{m/s}$。
  \item 角速度 $\omega$：单位时间转过的角度，$\mathrm{rad/s}$。
  \item 周期 $T$：转一圈的时间，$\mathrm{s}$。
\end{itemize}

更重要的是，高中揭示了一个关键结论：\textbf{匀速圆周运动需要指向圆心的合力}——这个力叫向心力。没有向心力，物体将沿切线飞出。

\section*{三、核心概念}

\begin{conceptbox}{描述圆周运动的物理量}

\textbf{线速度：}
\[
v = \frac{\Delta s}{\Delta t} = \frac{2\pi r}{T}
\]
方向沿圆周切线方向，时刻变化。

\textbf{角速度：}
\[
\omega = \frac{\Delta\theta}{\Delta t} = \frac{2\pi}{T}
\]
单位 $\mathrm{rad/s}$，方向用右手螺旋定则判断（高中阶段只关注大小）。

\textbf{周期与频率：}
\[
T = \frac{1}{f},\qquad \omega = 2\pi f
\]

\textbf{核心关系：}
\[
\boxed{v = \omega r}
\]

\begin{center}
% TikZ 图：匀速圆周运动
% Day 11：圆周运动
\begin{tikzpicture}[scale=0.8,>=Latex]
  \coordinate (P) at (30:2.5);
  % 轨道
  \draw[thick] (0,0) circle (2.5);
  % 圆心
  \filldraw[black] (0,0) circle (2pt);
  \node[below left] at (0,0) {$O$};
  % 半径与角度参考边
  \draw[dashed,gray] (0,0) -- (0:2.15);
  \draw[thick] (0,0) -- (30:2.5);
  \node[above left,fill=white,inner sep=1pt] at (48:1.05) {$r$};
  % 质点
  \filldraw[red] (P) circle (3pt);
  \node[above right] at (P) {$m$};
  % 速度方向（切线方向）
  \draw[->,thick,blue] (P) -- ++(120:1.2) node[above] {$\vect{v}$};
  % 向心加速度方向（指向圆心）：略微错开，避免遮住半径
  \draw[->,thick,red] ($(P)+(0.12,-0.18)$) -- ($(30:1.18)+(0.12,-0.18)$)
    node[pos=0.58,below right,fill=white,inner sep=1pt] {$\vect{a}_n$};
  % 角度
  \draw (0:0.8) arc (0:30:0.8);
  \node[fill=white,inner sep=1pt] at (15:1.0) {$\theta$};
\end{tikzpicture}


\figcaption{匀速圆周运动：$v$ 沿切线方向，$a_n$ 指向圆心}\label{fig:circular-motion}
\end{center}

\end{conceptbox}

\begin{conceptbox}{向心加速度}

\[
\boxed{a_n = \frac{v^2}{r} = \omega^2 r}
\]

\begin{itemize}
  \item 方向：始终指向圆心，与速度方向垂直。
  \item 物理意义：描述速度方向变化的快慢（不改变速度大小）。
  \item 即使是“匀速”圆周运动，也有加速度——因为速度方向不断变化，所以是\textbf{变加速曲线运动}。
\end{itemize}

\end{conceptbox}

\begin{conceptbox}{向心力}

\[
\boxed{F_n = m a_n = m\frac{v^2}{r} = m\omega^2 r}
\]

\begin{itemize}
  \item 向心力不是一种新的性质力，而是\textbf{指向圆心的合力}。
  \item 它由重力、弹力、摩擦力等真实力的合力（或其分力）提供。
  \item 匀速圆周运动中，合力完全提供向心力（大小恒定、方向始终指向圆心）。
\end{itemize}

\textbf{常见来源举例：}
\begin{itemize}
  \item 绳拉小球在水平面转圈：绳的拉力提供向心力。
  \item 汽车转弯：地面静摩擦力提供向心力。
  \item 天体绕行：万有引力提供向心力。
\end{itemize}

\end{conceptbox}

\section*{四、公式与适用条件}

\begin{center}
\begin{tabular}{lll}
\toprule
公式 & 物理意义 & 适用条件 \\
\midrule
$v = \omega r$ & 线速度与角速度关系 & 任何圆周运动 \\
$a_n = \dfrac{v^2}{r}$ & 向心加速度 & 匀速圆周运动（瞬时也适用于非匀速） \\
$F_n = m\dfrac{v^2}{r}$ & 向心力 & 合力指向圆心时 \\
$T = \dfrac{1}{f}$ & 周期与频率 & 周期性运动 \\
\bottomrule
\end{tabular}
\end{center}

\section*{五、典型模型}

\subsection*{模型 1：水平圆周运动（圆锥摆）}

轻绳一端固定，另一端系小球做水平圆周运动（圆锥摆）：
\[
T\sin\theta = m\omega^2 r = m\omega^2 l\sin\theta
\]
\[
T\cos\theta = mg
\]
解得 $\omega = \sqrt{\dfrac{g}{l\cos\theta}}$。

\subsection*{模型 2：竖直圆周运动（过山车）}

\begin{center}
% TikZ 图：竖直圆周运动（内侧轨道/过山车模型）
% Day 11：圆周运动
\begin{tikzpicture}[scale=0.72,>=Latex]
  \draw[thick] (0,0) circle (2.5);
  \filldraw[black] (0,0) circle (2pt);
  \node[left] at (0,0) {$O$};

  % 最高点：重力和支持力都指向圆心
  \filldraw[blue] (0,2.5) circle (3pt);
  \node[above] at (0,2.8) {最高点};
  \draw[->,thick,blue] (0,2.5) -- (1.05,2.5) node[above right] {$v$};
  \draw[->,thick,red] (-0.35,2.5) -- (-0.35,1.50) node[midway,left] {$mg$};
  \draw[->,thick,red] (0.35,2.5) -- (0.35,1.75) node[midway,right] {$N$};

  % 最低点：支持力向上，重力向下
  \filldraw[blue] (0,-2.5) circle (3pt);
  \node[below right] at (0.22,-2.82) {最低点};
  \draw[->,thick,blue] (0,-2.5) -- (1.05,-2.5) node[above right] {$v$};
  \draw[->,thick,red] (-0.35,-2.5) -- (-0.35,-3.35) node[midway,left] {$mg$};
  \draw[->,thick,red] (0.35,-2.5) -- (0.35,-1.55) node[midway,right] {$N$};

  \node[right] at (2.85,1.0) {最高点：$N+mg=m\dfrac{v^2}{r}$};
  \node[right] at (2.85,0.25) {临界：$N=0$};
  \node[right] at (2.85,-0.75) {最低点：$N-mg=m\dfrac{v^2}{r}$};
\end{tikzpicture}


\figcaption{竖直圆周运动最高点和最低点的受力分析}\label{fig:vertical-circular}
\end{center}

\textbf{最高点：} $N + mg = m\dfrac{v^2}{r}$，临界条件 $v_{\min} = \sqrt{gr}$（$N=0$ 时）。

\textbf{最低点：} $N - mg = m\dfrac{v^2}{r}$。

\subsection*{模型 3：汽车转弯}

水平路面：静摩擦力提供向心力，$f_s = m\dfrac{v^2}{r} \leq \mu_s mg$，最大安全速度 $v_{\max} = \sqrt{\mu_s g r}$。

\section*{六、例题精讲}

\begin{examplebox}{水平圆周——传动装置}

\textbf{题意分析：} 两个轮子用皮带传动，主动轮半径 $r_A = 20\ \mathrm{cm}$，从动轮半径 $r_B = 10\ \mathrm{cm}$。主动轮以 $\omega_A = 10\ \mathrm{rad/s}$ 转动。求：
\begin{enumerate}
  \item 从动轮的角速度 $\omega_B$。
  \item 两轮边缘点的线速度大小。
\end{enumerate}

\textbf{物理模型：} 皮带传动，两轮边缘线速度大小相等（皮带不打滑）。

\textbf{解题过程：}

\begin{enumerate}
  \item 线速度相等：$v_A = v_B$，$\omega_A r_A = \omega_B r_B$
  \[
  \omega_B = \frac{r_A}{r_B} \omega_A = \frac{20}{10} \times 10 = 20\ \mathrm{rad/s}
  \]
  
  \item $v = \omega_A r_A = 10 \times 0.2 = 2\ \mathrm{m/s}$
\end{enumerate}

\textbf{方法总结：} 皮带（或齿轮）传动中，边缘线速度相等；同轴转动中，角速度相等。

\end{examplebox}

\begin{examplebox}{汽车转弯}

\textbf{题意分析：} 汽车在水平弯道上转弯，弯道半径 $r = 50\ \mathrm{m}$，轮胎与地面间的最大静摩擦因数 $\mu_s = 0.4$，$g = 10\ \mathrm{m/s^2}$。求汽车安全转弯的最大速度。

\textbf{物理模型：} 静摩擦力提供向心力。

\textbf{解题过程：}

静摩擦力最大值 $f_{\max} = \mu_s mg$，提供向心力：
\[
\mu_s mg = m\frac{v^2}{r}
\]
\[
v_{\max} = \sqrt{\mu_s g r} = \sqrt{0.4 \times 10 \times 50} = \sqrt{200} \approx 14.1\ \mathrm{m/s}
\]

\textbf{方法总结：} 水平路面转弯的安全速度只与 $\mu_s$、$r$ 有关，与汽车质量无关。

\end{examplebox}

\begin{examplebox}{竖直圆周——过山车}

\textbf{题意分析：} 过山车在竖直圆形轨道上运行，轨道半径 $r = 8\ \mathrm{m}$，$g = 10\ \mathrm{m/s^2}$。求：
\begin{enumerate}
  \item 过山车在最高点不掉下来的最小速度。
  \item 若实际速度为 $v = 12\ \mathrm{m/s}$，求最高点时轨道对车的压力（$m = 500\ \mathrm{kg}$）。
\end{enumerate}

\textbf{物理模型：} 竖直圆周，最高点重力和支持力的合力提供向心力。

\textbf{解题过程：}

\begin{enumerate}
  \item 临界条件：支持力 $N = 0$，重力提供向心力：
  \[
  mg = m\frac{v^2}{r},\qquad v_{\min} = \sqrt{gr} = \sqrt{10 \times 8} = \sqrt{80} \approx 8.94\ \mathrm{m/s}
  \]
  
  \item $v = 12\ \mathrm{m/s} > v_{\min}$，安全：
  \[
  N + mg = m\frac{v^2}{r}
  \]
  \[
  N = m(\frac{v^2}{r} - g) = 500 \times (\frac{144}{8} - 10) = 500 \times (18 - 10) = 4000\ \mathrm{N}
  \]
\end{enumerate}

\textbf{方法总结：} 竖直圆周最高点的临界条件是 $N=0$、重力恰好提供向心力。速度小于临界值则会脱离轨道。

\end{examplebox}

\section*{七、变式训练}

\begin{enumerate}
  \item 一个 $m = 0.5\ \mathrm{kg}$ 的小球用 $l = 1\ \mathrm{m}$ 的细绳系着做圆锥摆，绳与竖直方向夹角 $\theta = 30^\circ$，$g = 10\ \mathrm{m/s^2}$。求小球运动的角速度和线速度。
  
  \item 汽车以 $v = 15\ \mathrm{m/s}$ 通过半径 $r = 30\ \mathrm{m}$ 的水平弯道，路面与轮胎间最大静摩擦因数 $\mu_s = 0.6$。问汽车是否侧滑？
  
  \item 质量为 $m$ 的小球在半径为 $r$ 的竖直光滑圆轨道内侧做圆周运动，在最高点速度为 $v$。求：
  \begin{enumerate}
    \item 轨道对小球的压力。
    \item 若 $v = \sqrt{2gr}$，求最高点压力。
  \end{enumerate}
\end{enumerate}

\section*{八、易错点}

\begin{mistakebox}
\begin{enumerate}
  \item \textbf{误以为向心力是“真实力”：} 向心力是效果力，不是性质力。受力分析时不能画“向心力”，只能画提供向心力的真实力（重力、弹力、摩擦力等）。
  
  \item \textbf{认为匀速圆周运动是匀速运动：} “匀速”只指速率不变，速度方向不断变化，所以有加速度。
  
  \item \textbf{竖直圆周中漏了重力：} 在最高点和最低点做受力分析时，重力始终存在。不要把 $N$ 当成唯一提供向心力的力。
  
  \item \textbf{传动问题中搞混传动方式：} 皮带传动 $\to$ 线速度相等；同轴转动 $\to$ 角速度相等。
  
  \item \textbf{汽车转弯误用滑动摩擦力：} 汽车正常转弯时是静摩擦力提供向心力，不是滑动摩擦力。临界打滑前可用 $f_{\max}=\mu_s mg$。
\end{enumerate}
\end{mistakebox}

\section*{九、竞赛入门补充}

\begin{competitionbox}
\textbf{非匀速圆周运动}

如果速率也在变化（如竖直圆周运动中重力沿切线方向的分力），向心加速度公式 $a_n = \dfrac{v^2}{r}$ 仍然成立，但此时还有切向加速度 $a_t$ 改变速度大小。

总加速度：
\[
a = \sqrt{a_n^2 + a_t^2}
\]

\textbf{例题：} 小球在竖直圆轨道最低点速度为 $v$，求该点的向心加速度和总加速度方向。

\textbf{解：} 在最低点，重力沿切线方向的分力为零，所以切向加速度 $a_t=0$。总加速度全部为向心加速度 $a_n = v^2/r$，方向竖直向上（指向圆心）。

\textbf{竞赛小贴士：} 变速圆周运动中，合力不指向圆心——它分解为指向圆心的法向分量（提供向心加速度）和沿切线方向的切向分量（提供切向加速度）。
\end{competitionbox}

\section*{十、今日小结}

\begin{enumerate}
  \item 线速度 $v = \omega r$，向心加速度 $a_n = \dfrac{v^2}{r} = \omega^2 r$。
  \item 向心力 $F_n = m\dfrac{v^2}{r}$，由真实力的合力提供，不是独立的性质力。
  \item 水平转弯：$f_s$ 提供向心力，$v_{\max} = \sqrt{\mu_s g r}$。
  \item 竖直圆周最高点临界：$v_{\min} = \sqrt{gr}$（$N=0$）。
  \item 皮带传动线速度相等，同轴转动角速度相等。
\end{enumerate}
